% !TITLE 今日良宴会
% !AUTHOR 无名氏
% !DYNASTY 两汉
% !GENRE 五言古诗
% !ORDER 4
% !SOURCE 《古诗十九首》；收于《文选》
% !EDITION 《古诗十九首集释》，隋树森，中华书局（2020年版为候选）
% !EDITION_STATUS pending

\begin{poem}
今日良宴会，欢乐难具陈\textsuperscript{1}。
弹筝奋逸响\textsuperscript{2}，新声妙入神\textsuperscript{3}。
令德唱高言\textsuperscript{4}，识曲听其真\textsuperscript{5}。
齐心同所愿，含意俱未申\textsuperscript{6}。
人生寄一世，奄忽若飙尘\textsuperscript{7}。
何不策高足，先据要路津\textsuperscript{8}？
无为守穷贱，轗轲长苦辛\textsuperscript{9}。
\end{poem}

\begin{annotations}
\notetext{1}{具陈：一一陈述。具，全部。}
\notetext{2}{奋逸响：发出奔放超逸的声音。奋，扬起。逸，超绝。}
\notetext{3}{新声：当时流行的新曲。妙入神：精妙到了出神的境地。}
\notetext{4}{令德：有美德的人，指座上德高望重者。唱高言：发出高妙的言论。唱，同“倡”，发出。}
\notetext{5}{识曲：懂得音乐的人。听其真：听出曲中的真意。}
\notetext{6}{齐心同所愿，含意俱未申：满座心意相同，却都藏在心里没有说出口。申，表达。}
\notetext{7}{奄忽：急速。飙尘：狂风卷起的尘土。}
\notetext{8}{策高足：鞭策骏马，比喻捷足先登。高足，良马。要路津：重要的路口和渡口，比喻显要的职位。}
\notetext{9}{无为：不要。轗轲（kǎn kē）：同“坎坷”，困顿不得志。}
\end{annotations}

\begin{appreciation}
《今日良宴会》是古诗十九首的第四首，也是这组诗里把欲望说得最直白的一首。

开头是一场热闹的宴会：“今日良宴会，欢乐难具陈”——今天的宴会真好，那份欢乐简直没法一一说尽。接着写音乐：“弹筝奋逸响，新声妙入神。”筝声扬起，奔放超逸；当时流行的新声，精妙得入了神。汉代人说“新声”，往往带着一点郑卫之音的意味，是市井里流行、士人私下喜欢的东西。

音乐之后是人言。“令德唱高言，识曲听其真”——座中有德行的长者发出高论，懂音乐的人听得出弦外之真意。这两句是过渡：琴声里有真意，人话里也有。真意是什么？“齐心同所愿，含意俱未申。”满座的人心里怀着同一个愿望，可谁也没有把它说出口。

这是全诗最耐人寻味的一处。一场宴会，人人都在笑，人人都在听，人人心里都揣着同一句话，却谁也不肯先开口。那句话是什么，下文替他们说了：“人生寄一世，奄忽若飙尘”——人寄居在这世上，快得像狂风里的一粒尘土。既然这样，“何不策高足，先据要路津？”为什么不鞭策骏马，抢先占住那要紧的路口和渡口？

“要路津”三个字，把整首诗的面目揭开了。所谓高言、真意，说到底是官位。路口和渡口是天下人必经之地，谁守在那里，谁就能收过往的利。诗人劝人趁早去占一个好位置，末了还补一句：“无为守穷贱，轗轲长苦辛。”不要守着穷困卑贱，一辈子坎坷辛苦。

儒家向来讲温柔敦厚，这一首却毫不遮掩。它不劝人修身，不劝人立德，只劝人做官，而且把做官说成抢路口。东汉后期，选举被世家大族把持，举孝廉、行征辟，名义上是选贤，实际上是攀附，当时已有童谣讽刺“举秀才，不知书；察孝廉，父别居”。一个没有门路的读书人，纵有满腹经纶，也只能“守穷贱”。所以这首诗的直白里有一种愤懑：既然规则如此，那就别谈什么清高了。
\end{appreciation}
